\documentclass[12pt]{article}
\usepackage[letterpaper,margin=1in]{geometry}
\usepackage{fontspec}
\usepackage{setspace}
\usepackage{amsmath}
\usepackage{amssymb}
\usepackage{microtype}
\usepackage{xcolor}
\usepackage{fancyhdr}
\usepackage{titlesec}
\usepackage{parskip}
\usepackage{hyperref}
\hypersetup{
  colorlinks=true,
  linkcolor=blue!50!black,
  urlcolor=blue!50!black,
  citecolor=blue!50!black,
  pdftitle={A1 to A2: The Actualization of Perception (Synthesized)},
  pdfauthor={Dalton Salvo},
  pdfsubject={Dissertation -- Aristotelian Actualization Chain}
}

% Main font: Linux Libertine if available, otherwise default serif
\IfFontExistsTF{Linux Libertine O}{%
  \setmainfont{Linux Libertine O}[Numbers=OldStyle]%
}{%
  \IfFontExistsTF{TeX Gyre Pagella}{%
    \setmainfont{TeX Gyre Pagella}%
  }{}%
}

% Greek wrapping: try dedicated Greek font, else fall back to main font
\IfFontExistsTF{GFS Didot}{%
  \newfontfamily{\greekfont}{GFS Didot}%
  \newcommand{\gk}[1]{{\greekfont #1}}%
}{%
  \IfFontExistsTF{DejaVu Serif}{%
    \newfontfamily{\greekfont}{DejaVu Serif}%
    \newcommand{\gk}[1]{{\greekfont #1}}%
  }{%
    \newcommand{\gk}[1]{#1}%
  }
}

% Line spacing and paragraph style
\setstretch{1.3}
\setlength{\parindent}{0pt}
\setlength{\parskip}{0.8em}

% Section titles
\titleformat{\section}{\Large\bfseries}{}{0pt}{}
\titleformat{\subsection}{\large\bfseries}{}{0pt}{}
\titlespacing*{\section}{0pt}{2em}{1em}
\titlespacing*{\subsection}{0pt}{1.5em}{0.6em}

% Footnote style
\usepackage[hang,flushmargin]{footmisc}

% Page header/footer
\setlength{\headheight}{14pt}
\pagestyle{fancy}
\fancyhf{}
\fancyhead[L]{\small\itshape A\textsubscript{1}$\to$A\textsubscript{2}: Aisthēsis (Synthesis)}
\fancyhead[R]{\small\itshape Draft --- \today}
\fancyfoot[C]{\thepage}
\renewcommand{\headrulewidth}{0.4pt}
\renewcommand{\footrulewidth}{0pt}

% Title page content
\title{\vspace{-2em}\textbf{$A_1 \to A_2$: The Actualization of Perception}\\[0.5em]
       \large\normalfont A Dissertation Section (Synthesized Draft)}
\author{Dalton Salvo}
\date{April 2026 --- Draft for Review}

\begin{document}
\maketitle
\thispagestyle{fancy}

\subsection*{1. $A_1 \to A_2$: The Actualization of Perception}

The preceding section established that motion and time constitute $A_0$, the ground-level actuality from which every subsequent stage of the perceptual chain arises. Motion, understood as the fulfilment of what is potentially insofar as it is potential, provides the ontological scaffold upon which the ensouled body first encounters the sensible world; time, as the number of motion with respect to the before and after, supplies the continuous metric within which successive actualizations can be ordered and distinguished (Aristotle, \textit{Physics}, p.~60). We turn now to the passage from $A_1$ to $A_2$---the actualization of perception itself, \textit{aisthēsis}---which names the specific \textit{kinēsis} whose unmoved originator is the sensible object exercising its proper quality in actuality, whose moved mover is the sense faculty receiving the sensible form without the matter, and whose terminus is the completed perceptual act together with the dual residual trace it deposits in the ensouled body. This transition is perhaps the decisive hinge of the entire actualization chain, for it is here that the physical world first acquires phenomenal determinacy for the living creature, and here that the material conditions for \textit{phantasia} and all subsequent cognitive operations are first laid down. Aristotle frames the soul as ``a substance in the sense of the form of a natural body'' (Aristotle, \textit{On The Soul (De Anima)}, p.~18), thereby situating the perceptual faculty within the hylomorphic constitution of the organism rather than treating it as a detached cognitive module. The sense faculty, in its first actuality, is a capacity---a \textit{hexis} or structured readiness---that awaits actualization by an external agent. The sensible object, for its part, possesses its quality in full actuality: the red of the rose, the warmth of the flame, the pitch of the string. When these two first actualities are brought into the appropriate spatial and temporal relation through a properly disposed medium, the transition from $A_1$ to $A_2$ occurs.

The section proceeds through seven architectural moves whose joint effect is to specify the perceptual chain from preconditions to deposited residue. First, it specifies the $A_1$ preconditions and the medium that transmits the sensible form. Second, it develops Aristotle's hylomorphic account of form-reception through the wax-signet metaphor of \textit{De Anima} II.12, 424a17--24. Third, it applies the three-factor kinetic schema established at $A_0$ to perception, showing that the activity of the sensible object and the activity of the sense faculty are numerically one event. Fourth, it analyses perception as an enmattered \textit{alloiōsis} of the preservative kind rather than the corrupting kind. Fifth, it establishes that the perceptual motion does not terminate with the object's withdrawal but persists as residual motion in the sensory apparatus---what this chapter names resonant \textit{kinēsis}. Sixth, it advances the dual-trace thesis: the residual motion is not a single undifferentiated impression but a composite whose formal-epistemic aspect (resonant \textit{aisthēma}) and affective-orectic aspect (resonant \textit{orexis}) are one in substrate and two in being. Seventh, it situates $A_2$'s dual-trace composite as the unmoved originator of the phantastic motion $M_2 \to A_3$. A Development Note flags open questions that subsequent sections will take up.

The architectural stakes are considerable. If $A_2$'s completed actuality harboured only a formal trace, the chain could not explain how perception generates the motivational states that drive animal action. The dual-trace thesis developed in §7 is thus not an ornament but a structural requirement: it specifies the material from which the chain's subsequent stages---the \textit{phantasma} at $A_3$, the cognitive actualities and the committal dimension of \textit{doxa} at $M_2 \to M_3$, the determinate \textit{pathē} that mobilize \textit{orexis} at $M_3 \to M_4$---will be generated. Aristotle does not name the second trace; the reading advanced here treats the dual structure as a structural requirement of his own commitments, to be defended by close engagement with the texts that ground it.

Burke's analysis of motivation in terms of action and substance provides an instructive parallel here, insofar as he insists that terminologies of motion and conditioning are dialectical enterprises designed to clarify the relationship between agent and scene (Burke, \textit{A Grammar of Motives}, pp.~77--80). Similarly, the Aristotelian account of perception is not merely a causal story about physical impingement but a dialectical articulation of how the scene---the sensible environment---and the agent---the perceiving organism---are mutually implicated in a single actualization. What distinguishes the perceptual motion from ordinary physical alteration is precisely its hylomorphic character: the sense faculty receives the form of the sensible object without receiving its matter. This formulation, which has generated extensive scholarly debate, establishes perception as a unique species of \textit{alloiōsis}, one in which the organ is altered not by taking on the material substrate of the quality but by being assimilated to its formal structure. Gonzalez, examining Aristotle's account of persuasive speech in the \textit{Rhetoric}, emphasizes the reciprocal character of such formal reception, observing that ``the orator succeeds when his \textit{eunoia} and \textit{philia} are reciprocated by his audience'' (Gonzalez, \textit{The Meaning and Function of Phantasia in Aristotle's `Rhetoric' III.1}, p.~29). Though Gonzalez's immediate concern is rhetorical rather than perceptual, the structural parallel is instructive: just as the orator's goodwill must be received and reciprocated by the audience for persuasion to occur, so too must the sensible form be received and assimilated by the sense organ for perception to be actualized. In both cases, a dyadic relation of active and receptive poles must be brought to completion in a single shared actuality. Gibson, working within a radically different theoretical framework, nonetheless arrives at a convergent insight when he argues that ``motion and rest are in fact what an observer experiences with flow and nonflow of the array'' (Gibson, \textit{The Ecological Approach to Visual Perception}, pp.~209--210), thereby insisting that perception is not a passive registration of discrete stimuli but an active engagement with the structured affordances of the environment. The sections that follow will trace each moment of this actualization chain in detail, beginning with the preconditions that must obtain before the perceptual act can occur.

\subsection*{2. The Preconditions: $A_1$ and the Medium}

Before the perceptual act can be actualized, two distinct first actualities must be in place: the sensible object must be exercising its proper quality, and the sense organ must possess its capacity in an appropriately disposed state. These twin preconditions constitute $A_1$, the stage of readiness that precedes but does not yet amount to perception. Aristotle's treatment of this stage is distributed across several passages of the \textit{De Anima}, but its logic is perhaps clearest in his account of the distinction between first and second actuality. Aristotle develops this distinction most explicitly at \textit{De Anima} II.5, 417a21--b2: the geometer who holds geometrical knowledge without currently exercising it possesses a first actuality---a \textit{hexis}, a stable disposition capable of being exercised---while the geometer who actively proves a theorem exercises that capacity and thereby possesses a second actuality. The soul, as the first actuality of a natural body potentially having life, stands to the body as the sailor's knowledge of navigation stands to the act of navigating; the capacity is real and fully constituted, yet it awaits an occasion and an object to be exercised (Aristotle, \textit{On The Soul (De Anima)}, p.~18). Accordingly, the sense faculty in its first-actual state is not merely a bare potentiality---it is not the potentiality of an infant who has never perceived---but a structured readiness, a \textit{hexis}, that has been informed by the organism's developmental history and is now poised for activation. This first actuality is not a bare potentiality to be awakened from without; it is already an actuality, a structured readiness that awaits only the proper object to pass into the second actuality of exercise.

The sensible object, for its part, must not merely possess its quality dispositionally but must be actively exercising it. Aristotle distinguishes carefully between what is ``capable of being seen'' and what is ``being seen'' (Aristotle, \textit{Metaphysics}, pp.~105--106); the former names the first actuality of the visible quality, while the latter names the second actuality in which that quality is fully present to a perceiver. A coloured surface in a lightless room possesses the first actuality of visibility without exercising it; only when illuminated and placed within the perceptual field of a suitably disposed organism does the quality pass from dispositional to active presence. This distinction is crucial because it establishes that perception requires a genuine exercise of causal power on the part of the object; the sensible quality is the unmoved mover of the perceptual chain, originating the motion without itself undergoing change. Aristotle does not develop a separate technical analysis of sensible qualities in act, because the point is already available from the $A_0$ framework: the sensible object is an \textit{entelecheia} whose \textit{telos} is internal, and its exercise requires no further enabling condition beyond the motion and time in which all natural activity occurs.

Between object and organ, however, lies the medium---the diaphanous in the case of sight, the air or water in the case of hearing and smell. The medium is not merely an inert conduit through which the sensible form travels; it is itself a necessary precondition, a structural feature of the perceptual situation that must be properly disposed. Aristotle insists that direct contact between object and organ does not produce perception in the case of the distance senses; there must be an intervening medium through which the form is transmitted. He develops the medium's role most explicitly at \textit{Sense and Sensibilia} 6, 447a3--7: ``if the body which is heated or frozen is extensive, each part of it successively is affected by the part contiguous, while the part first changed in quality is so changed by the cause itself which originates the change.'' The medium's affection is sequential. The sensible quality alters the part of the medium in contact with the object; that altered part alters the next part; and so on until the alteration reaches the sense-organ. Sound, whose propagation Aristotle analyses most explicitly, is genuinely successive: we perceive the clap of thunder after we see the flash because the alteration of the air takes measurable time to traverse the medium. Light is the exception. Aristotle notes at \textit{Sense and Sensibilia} 6, 446b8--9 that light is transmitted instantaneously---not propagated successively through the medium but actualizing the transparent in a single movement. The sequential character characteristic of sound and heat does not obtain for vision.

Burke, analyzing the scenic dimensions of motivated action, argues that the ``circumference of the scene in which occurs the given act'' can never be fully delimited, since an act always has a background that extends beyond any particular boundary (Burke, \textit{A Grammar of Motives}, pp.~383--384). The Aristotelian medium functions as precisely such a scenic circumference---an enabling condition that is neither agent nor patient but the field within which the perceptual transaction occurs. Furthermore, as Aristotle notes in his account of incidental perception, ``we speak of an incidental object of sense where e.g. the white object which we see is the son of Diares; here because being the son of Diares is incidental to the white which is perceived, we speak of the son of Diares as being incidentally perceived'' (Aristotle, \textit{On The Soul (De Anima)}, pp.~25--27). This passage illuminates the layered character of $A_1$: the proper sensible---whiteness---is what the sense faculty is directly poised to receive, but the total perceptual situation includes incidental features that are co-present without being directly actualized by the sense organ's specific capacity.

Gibson's ecological approach offers a modern reformulation of these Aristotelian preconditions. Where Aristotle speaks of the medium as diaphanous---transparent and capable of receiving and transmitting the form of colour---Gibson speaks of the ambient optic array as a structured field of light that specifies the surfaces and layouts of the environment (Gibson, \textit{The Ecological Approach to Visual Perception}, pp.~209--210). In both accounts, perception is not a matter of bare stimulus impinging upon a passive receptor but rather a structured encounter between an already organized environment and an already capacitated organism. The medium, whether conceived as Aristotelian diaphanous or as Gibsonian ambient array, is the enabling third term that makes the encounter possible. The medium's sequential affection is, finally, the physical correlate of the temporal horizon $A_0$ establishes: where $A_0$ supplies the general ontological conditions---motion and time as co-actual---the medium supplies the specific physical mechanism through which one particular species of motion, the perceptual kind, unfolds in time. Thus $A_1$ names not a single state but a constellation of preconditions---object, medium, and organ---each of which must be in its proper first actuality before the transition to $A_2$ can commence; and each of which is itself made possible by the motion and time already in act at $A_0$.

\subsection*{3. Hylomorphism and the Reception of Form}

The central and most controversial claim of Aristotle's perceptual theory is that the sense faculty receives the sensible form without the matter. This formulation, which has occasioned a vast literature of interpretation, is the hylomorphic core of the $A_1 \to A_2$ transition: it specifies the precise manner in which the sense organ is altered by its encounter with the sensible object. The governing text is \textit{De Anima} II.12, 424a17--24: ``a sense is what has the power of receiving into itself the sensible forms of things without the matter, in the way in which a piece of wax takes on the impress of a signet-ring without the iron or gold; what produces the impression is a signet of bronze or gold, but not \textit{qua} bronze or gold: in a similar way the sense is affected by what is coloured or flavoured or sounding not insofar as each is what it is, but insofar as it is of such and such a sort and according to its form.'' To receive the form without the matter is to undergo a change that is formally but not materially identical with the quality of the object; the eye, in seeing red, does not become materially red in the way that a dye stains cloth, but it does take on the formal structure of redness, becoming, as it were, isomorphic with the quality it perceives. The signet-ring analogy is architecturally decisive. The ring impresses its engraved pattern into the wax; the wax takes the pattern; but the wax does not become iron or gold. What passes from ring to wax is the \textit{form}---the structured pattern of convexities and concavities---not the substrate that bears the form. Perception operates analogously: the colour of the apple passes to the eye, the pitch of the string passes to the ear, but the eye does not become apple-matter, nor the ear string-matter.

Aristotle distinguishes the organ as spatial magnitude from the sensory power as form-in-magnitude at \textit{De Anima} II.12, 424a24--29: ``The sense and its organ are the same in fact, but their essence is not the same. What perceives is, of course, a spatial magnitude, but we must not admit that either the having the power to perceive or the sense itself is a magnitude; what they are is a certain form or power in a magnitude.'' The distinction is crucial. The eye is a spatial magnitude---a physical object with extension, composed of matter, located in the body. The \textit{capacity} to see, however, is not itself a magnitude; it is a form present in the magnitude of the eye, a ratio or structured configuration that the organ's material substrate realizes. When the organ receives the sensible form, it is the form-in-magnitude that is altered, not the material magnitude. The material magnitude is unchanged in kind; what changes is the formal configuration that the material magnitude realizes.

The form-in-magnitude that is the sensory power, Aristotle argues, is a \textit{logos}---a ratio of the contraries proper to the sense's domain. At \textit{De Anima} II.11, 424a4--5 he observes that sense is ``a sort of mean'' between the extremes of the sensible qualities, and at \textit{De Anima} III.2, 426a27--b8 he develops the ratio analysis explicitly. Each sense is constituted by a proportion: the eye between light and dark, the ear between high and low pitch, the tongue between sweet and bitter. The \textit{logos} is structurally a ratio, and perceptible qualities are themselves ratios of their underlying contraries---grey as a specific proportion of white and black, a flavour as a proportion of sweet and bitter. Perception, on this account, is the organ's taking-on of the ratio that the perceptible quality instantiates: the sensory power adjusts its ratio to match the object's ratio, and this formal matching \textit{is} the perceptual act.

The \textit{logos} framework explains the destruction condition that Aristotle develops at \textit{De Anima} III.2, 426b3--8: if the sensible exceeds the ratio of the organ, the motion acts on the organ \textit{as matter} rather than \textit{as sense}. Excess brightness destroys sight; excess loudness destroys hearing; excess heat destroys touch. When the stimulus intensity exceeds the proportional range the organ can receive, the organ is affected as matter rather than as sense, and the \textit{logos} itself is destroyed rather than realized. The motion produced is real, and it is genuine damage to the organ, but it is not perception. Perception requires that the stimulus fall within the ratio the organ can receive; excess falls outside the ratio, and the organ suffers rather than perceives. This destruction condition presupposes the hylomorphic analysis: if perception were simply material alteration, there would be no principled difference between perceiving and being damaged. The \textit{logos} is what makes that difference principled.

The hylomorphic account of perception thus distinguishes sharply between ordinary physical alteration---in which one body transmits both form and matter to another, as fire heats and dries wood---and perceptual alteration, in which form alone is transmitted. Nussbaum, in her analysis of \textit{phantasia}'s role in Aristotle's account of action, touches upon the significance of this distinction when she observes that ``\textit{phantasmata} are like \textit{aisthēmata}, only `without matter''' (Nussbaum, \textit{The Role of Phantasia in Aristotle's Explanation of Action}, p.~47). If perception did not involve the reception of form without matter, there would be no basis for the subsequent claim that \textit{phantasia} preserves formal structures in the absence of their original material substrates; hence the hylomorphic theory of perception is not merely a local doctrine about sense organs but the foundation of Aristotle's entire cognitive psychology. White, situating this doctrine within the broader architectonic of the \textit{De Anima}, identifies the soul's critical power---\textit{to kritikon}---as the unifying thread that connects perception, \textit{phantasia}, and thought (White, \textit{The Meaning of Phantasia in Aristotle's De Anima III, 3-8}, pp.~9--11). The soul's capacity to discriminate---to judge that this is white and not black, sweet and not bitter---presupposes that it has received the form of the sensible quality in a manner sufficiently determinate to ground such judgment. The hylomorphic reception of form is thus not a mere passive imprinting but a discriminative event; the sense faculty, in receiving the form, simultaneously exercises its critical capacity to distinguish what it has received.

Frede, examining the cognitive role of \textit{phantasia}, notes that \textit{phantasiai} have sometimes been ``degraded to mere epiphenomena, the lingering after-images of sensations'' (Frede, \textit{The Cognitive Role of Phantasia in Aristotle}, p.~3). This characterization, however, rests upon a reductive understanding of the hylomorphic reception that underlies perception: if we take the reception of form without matter seriously, then the formal structure deposited in the sense organ during perception is not an epiphenomenal residue but a genuine ontological determination of the ensouled body. The persistence of this formal determination after the external object is removed is what makes \textit{phantasia} possible, and it is the hylomorphic character of the original reception that ensures that the persisting trace preserves the formal structure of the perceived quality. The wax-signet analogy accordingly does more than illustrate form-reception; it marks the specific mode of being that perception exhibits. The wax undergoes genuine alteration---it acquires a pattern it did not previously have---but its material substrate is preserved rather than transformed. The alteration is formal, not material. This points forward to the next architectural move: perception is an \textit{alloiōsis}, but not every \textit{alloiōsis} is corruption, and the distinction between two kinds of qualitative alteration, and the assignment of perception to the second kind, must now be developed.

\subsection*{4. The Activity of Mover and Moved as One}

Aristotle's account of the perceptual act turns upon a principle drawn from the \textit{Physics}: the activity of the mover and the activity of the moved are one and the same activity, though they differ in their being. The three-factor kinetic schema established at $A_0$---unmoved originator, moved mover, that which is moved---applies directly to perception. The unmoved originator at $A_1 \to A_2$ is the sensible object in actuality, whose form initiates the motion without being itself altered by it. The moved mover is the sense faculty, which is altered (formally, per §3) in the very act of receiving the sensible form and which thereby moves the perceiver as a whole. That which is moved is the perceiver via the sense-organ. The schema iterates from $A_0$ with the same structural logic: one actuality functions as the unmoved source of the motion that generates the next actuality.

Decisive for the perceptual application is Aristotle's claim at \textit{Physics} III.2, 202a13--20 that the actuality of the mover and the actuality of the moved are numerically one. ``Motion is in the movable. It is the fulfilment of this potentiality by the action of that which has the power of causing motion; and the actuality of that which has the power of causing motion is not other than the actuality of the movable; for it must be the fulfilment of \textit{both}.'' One motion, not two. The mover's \textit{poiesis} and the moved's \textit{pathesis} are distinguished in account but not in substrate; they are one event described under two categorial aspects. Aristotle makes the point concrete with the Thebes--Athens analogy at \textit{Physics} III.3, 202b13--14: ``the road from Thebes to Athens and the road from Athens to Thebes are the same.'' One road, two descriptions. One motion, two formal aspects.

\textit{De Anima} III.2, 425b26--27 extends this directly to perception: ``the activity of the sensible object and that of the sense is one and the same activity, and yet the distinction between their being remains.'' The sensible object's exercise of its power (its \textit{poiesis}) and the sense faculty's reception of the form (its \textit{pathesis}) are not two events conjoined but one event under two \textit{logoi}. Aristotle illustrates with a striking example at \textit{De Anima} III.2, 425b28--426a1: ``a man may have hearing and yet not be hearing, and that which has a sound is not always sounding. But when that which can hear is actively hearing and that which can sound is sounding, then the actual hearing and the actual sound come about at the same time.'' Hearing and sounding---the active hearing and the active sounding---arise together, in a single actualization. The sound is not first produced and then heard; the sounding and the hearing are the same actuality under two descriptions.

The formula Aristotle employs here---one in substrate, different in being---will recur at every subsequent stage of the chain. It is introduced at the motion-level in \textit{Physics} III.2--3 (mover and moved, \textit{poiesis} and \textit{pathesis}); it is redeployed for perception at \textit{De Anima} III.2 (sensible and sense); it will be applied again at \textit{De Anima} III.7, 431a13--14 to the coordination of perception and appetite (one in substrate, different in being); and it structures the dual-trace thesis to be developed in §7. The formula's persistence across contexts reflects the chain's architectural unity: a single ontological pattern---numerical identity under plural formal descriptions---iterates across levels, each time specifying a new relation that the substrate-level unity secures.

This doctrine has profound consequences for the topology of the perceptual event. If the act of the sensible object and the act of the sense organ are one, then perception cannot be adequately described as a chain of discrete causal steps---stimulus, transmission, reception---but must be understood as a single unified event that spans the distance between object, medium, and organ. Burke, analyzing the rhetorical situation in analogous terms, observes that a work's fitness---its decorum---involves a purely technical adaptation of artistic means to artistic end that simultaneously implicates social propriety (Burke, \textit{A Rhetoric of Motives}, p.~149). Similarly, the perceptual act is a technical adaptation of organ to object that simultaneously implicates the broader totality of involvements within which the organism exists. The act is one, but its description requires distinguishing the contributions of agent and patient, just as the description of a rhetorical performance requires distinguishing speaker, speech, and audience even though the persuasive event is a single occurrence. Gonzalez, examining the function of \textit{phantasia} in rhetoric, underscores this relational character by noting that effective speech depends upon a reciprocity between orator and audience: ``the orator succeeds when his \textit{eunoia} and \textit{philia} are reciprocated by his audience'' (Gonzalez, \textit{The Meaning and Function of Phantasia in Aristotle's `Rhetoric' III.1}, p.~29). The perceptual analogue is precise: perception succeeds when the sensible object's formal power is reciprocated by the sense organ's receptive capacity. Without this reciprocity, no actualization occurs---the form remains unperceived, the capacity unexercised. The unity of mover and moved in the perceptual act thus establishes a fundamental principle of correlation: the sensible world and the sentient organism are not independent substances that happen to interact but correlative poles of a single actualization, each requiring the other for its full reality to be disclosed.

One further feature of the perceptual motion must be registered. Because the mover's actuality and the moved's actuality are numerically one, the perceptual event does not consist in the mover producing something that then exists separately from the mover in the moved. The actuality \textit{is} the shared event, and the moved's being-affected is an aspect of that event rather than a subsequent state it acquires. This structural feature licenses what §6 will develop: the event, once complete, does not simply vanish. Motion that is the shared actuality of two substances leaves traces in the moved that are the residual character of that shared actuality. Heidegger's later analysis of \textit{Befindlichkeit}---the attunement or dispositional state in which Dasein always already finds itself---recasts this Aristotelian insight in existential terms; just as the sense organ must already be disposed to receive the form of the sensible object, so too must Dasein be attuned to a world that has already disclosed itself in a particular mood (Multiple Authors, \textit{Heidegger and Rhetoric}, pp.~113--115). In both cases, receptivity is not passivity but a structured readiness that is constitutive of the perceptual or existential situation.

\subsection*{5. Perception as Enmattered Alloiōsis}

Although the reception of form without matter is the defining characteristic of perception, Aristotle is careful to insist that this reception is not a purely formal event detached from the body. Perception is an enmattered \textit{alloiōsis}---a qualitative alteration that occurs in and through the material constitution of the sense organ. Aristotle frames perception as a species of \textit{alloiōsis} explicitly at \textit{De Anima} II.5, 416b33--417a1: ``Sensation depends, as we have said, on a process of movement or affection from without, for it is held to be some sort of change in quality.'' The eye jelly is physically affected by light; the tympanic membrane vibrates in response to sound; the flesh is heated or cooled by the tangible object. These physiological changes are not mere accompaniments of perception but its material substrate; the formal reception could not occur without them. Aristotle's own formulation is unambiguous: the soul is ``a substance in the sense of the form of a natural body'' that has life potentially (Aristotle, \textit{On The Soul (De Anima)}, p.~18), and the activities of the soul---including perception---are inseparable from the body whose form it is.

Two kinds of \textit{alloiōsis} must be distinguished. The first is ordinary qualitative change, in which a substance's quality is altered by the imposition of a contrary: cold water becomes hot, a living body becomes diseased, a flame is extinguished. This kind of \textit{alloiōsis} is corrupting in the sense that it destroys the subject's prior state; the contrary supplants what was there. The second kind of \textit{alloiōsis}---which Aristotle develops at \textit{De Anima} II.5, 417a21--b2---is preservative rather than corrupting. ``When we use the term `to be acted on,' we do not use it in one single sense, but in two: one is the destruction of one of two contraries by the other, the other is the preservation of what is potential by what is actual and like it.'' The geometer learning geometry, the builder acquiring the art of building, the sense-organ receiving a sensible form: these are alterations that do not destroy the subject's prior state but actualize a potentiality that was already present as first actuality. The preservative \textit{alloiōsis} is not a process in which the subject is replaced by its contrary; it is a process in which a capacity already present in the subject is brought to exercise.

Perception belongs to the preservative kind. When the eye sees red, the eye is not thereby destroyed or replaced by red; its capacity to see is actualized, and the actualization is a \textit{fulfilment} of that capacity rather than its negation. Aristotle is explicit at \textit{De Anima} III.7, 431a3--5: ``in the case of sense clearly the sensitive faculty already was potentially what the object makes it to be actually; the faculty is not affected or altered'' in the corruptive sense. The sensitive faculty receives the sensible form by actualizing what it was already capable of being; the reception does not introduce a foreign element that destroys the organ's prior state. The wax-signet analogy clarifies the point. The wax is altered by the ring---it acquires a pattern it did not previously have---but the alteration is not the wax's corruption; the wax is preserved as wax, and the pattern it receives is a determination of the wax rather than a substitution for it. The sensory organ is likewise preserved as what it is while receiving the form of what it perceives; the reception is a determination of the organ, not a substitution.

Thus perception as enmattered \textit{alloiōsis} names the inseparability of formal and material dimensions within the perceptual act. Nussbaum draws attention to this enmattered character when she argues that the activities of \textit{nous} in a living body ``must be linked with those of the faculty of \textit{phantasia}, which it would not be implausible to see as realized in some physiological occurrences'' (Nussbaum, \textit{The Role of Phantasia in Aristotle's Explanation of Action}, pp.~26--27). If even noetic activity requires a physiological realization through \textit{phantasia}, then \textit{a fortiori} perception, which is the causal origin of \textit{phantasia}, must be thoroughly enmattered. The physiological changes that accompany perception---the warming of the organ, the reception of a particular ratio of qualities---are not accidental but essential features of the perceptual process. Accordingly, the $A_1 \to A_2$ transition is not a passage from matter to form but a passage within enmattered form; the organ, already constituted as a determinate material structure with a specific formal capacity, undergoes a further formal determination that alters its material state without destroying its underlying constitution.

This conception of perception as enmattered \textit{alloiōsis} has important consequences for the understanding of the residual trace. If perception involved only the formal reception of a quality without any material alteration, it would be difficult to explain how the sensible form persists in the organ after the external object is removed. But because perception is an enmattered process---because the formal reception is accompanied by and realized in a physiological change---the persistence of the trace is explicable as a material residue of the perceptual event. The physiological alteration persists, and with it the formal structure that it realizes. White's identification of \textit{to kritikon} as the unifying capacity of the discriminative soul illuminates this point: the soul's power to discriminate---to distinguish white from black, sweet from bitter---is exercised through and in the physiological processes of the sense organs (White, \textit{The Meaning of Phantasia in Aristotle's De Anima III, 3-8}, pp.~9--11). Moreover, Burke's analysis of how ``body'' in even its ``lowliest'' forms can be treated as ``a way into `spirit''' rather than as its antithesis provides a conceptual frame for understanding the Aristotelian refusal to separate material and formal dimensions of perception (Burke, \textit{A Rhetoric of Motives}, p.~203). In both frameworks---the Aristotelian and the Burkean---matter and form, body and soul, are not opposed terms but co-constitutive dimensions of a single process. Gibson's ecological psychology similarly resists the separation of perceiver and environment, insisting that the observer's experience of motion and rest is constituted by the optical flow of the ambient array rather than by discrete sensory inputs processed by a disembodied mind (Gibson, \textit{The Ecological Approach to Visual Perception}, pp.~209--210). The doctrine that ``optical motion perspective is not the same as visual kinesthesis'' implies that what the perceiver encounters is not a representation of the world but the world itself as specified by the structured energy available to the organism. This convergence between Aristotelian enmattered \textit{alloiōsis} and Gibsonian ecological perception reinforces the claim that the $A_1 \to A_2$ transition is an event that occurs at the intersection of organism and environment, form and matter, rather than in a disembodied cognitive space.

The structural consequence matters decisively for what follows. If perception were a corruptive \textit{alloiōsis}---if the organ were destroyed or substantially transformed in each perceptual act---then nothing would remain of the act once it ended. The organ would have to be restored before the next perception could occur. But perception is not corruptive; it is preservative, and the preservation means that the form received leaves a trace. The organ, having been formally altered, does not return to its pre-perceptual state the moment the object withdraws. Something persists. What persists, and how it does so, is the subject of the next subsection.

\subsection*{6. Resonant Kinēsis: The Persisting Impression}

The completion of the perceptual act does not mark the termination of the \textit{kinēsis} initiated by the sensible object; rather, it deposits a persisting impression in the ensouled body, a residual movement that continues after the object is no longer present. Aristotle makes this point at several places in the corpus with striking consistency. At \textit{De Anima} III.2, 425b23--26 he observes: ``even when the sensible objects are gone the sensings and imaginings continue to exist in the sense-organs.'' At \textit{On Dreams} 2, 459a24--28 he elaborates: ``The objects of sense-perception corresponding to each sensory organ produce sense-perception in us, and the affection due to their operation is present in the organs of sense not only when the perceptions are actualized, but even when they have departed.'' The residual affection is not a metaphor for an after-effect; it is the continuation of the perceptual motion in the organ after the originating cause has withdrawn. The nature and status of this residual \textit{kinēsis} are thus of the highest importance for the actualization chain, for it is through this impression that the completed perceptual act ($A_2$) becomes the precondition for the next stage of actualization ($A_3$).

Aristotle illustrates the mechanism with a physical analogy at \textit{On Dreams} 2, 459b1--6: ``This we must likewise assume to happen in the case of qualitative change; for that part which has been heated by something hot, heats the part next to it, and this propagates the affection onwards to the starting-point. This must therefore happen in sense-perception, since actual perceiving is a qualitative change.'' The analogy compares the residual motion to heat remaining in water after the fire is removed. The water, having been heated, does not instantaneously return to its pre-heated state when the fire is withdrawn; the warmth persists as a residual motion in the water's material substrate, gradually dissipating as the water equilibrates with its surroundings. So too in perception: the organ, having been formally altered, retains the alteration as a residual motion that persists for some time after the object has withdrawn. Aristotle gives a name to these residual affections at \textit{On Dreams} 2, 460b2--3: ``even when the external object of perception has departed, the impressions it has made persist, and are themselves objects of perception.'' The Greek is \gk{αἰσθήματα}---\textit{aisthēmata}, impressions or residual sensings. The \textit{aisthēmata} are the persisting formal traces of the completed perceptual act; they are ``themselves objects of perception'' in the sense that they can be attended to by the sensory apparatus in the original object's absence; they are the material of dreams, afterimages, imagination, and the beginnings of memory.

This chapter adopts the term \textit{resonant kinēsis} for the residual motion understood as a whole---the ongoing activity in the sensory apparatus that persists after the sensible object has withdrawn. The coinage marks two features of the trace that Aristotle's own lexicon does not consolidate into a single technical vocabulary. First, the trace is \textit{kinēsis}---it is a motion, not a static imprint, and it has the ongoing character of all Aristotelian \textit{kinēsis}. Second, it is \textit{resonant}---it continues reverberating within the organ like sound continuing in a stilled bell, retaining the formal structure of the original motion while operating in the absence of its originating cause. Kevin White captures the reverberating character precisely: ``This lingering, resonating, echoing presence of sensible forms freed from their original matter appears to be the essence of \textit{phantasia} for Aristotle. The movement in the sense-power caused by the sensible object itself sets up a second movement which continues after the sensation has ceased and the object is gone'' (White, \textit{The Meaning of Phantasia in Aristotle's De Anima III, 3-8}, p.~498). The residual motion is not a weakened copy of the original; it is a continuing exercise of the same formal structure, sustained by the organ's material substrate in the original cause's absence.

Frede, examining the cognitive role of \textit{phantasia}, notes that \textit{phantasiai} have often been treated as ``the lingering after-images of sensations'' (Frede, \textit{The Cognitive Role of Phantasia in Aristotle}, p.~3). However, as Frede notes, this account does not do justice to the full range of functions that \textit{phantasia} performs; it reduces the persisting trace to a degraded copy of the original perception rather than recognizing it as a formal determination in its own right (Frede, \textit{The Cognitive Role of Phantasia in Aristotle}, pp.~4--5). The resonant \textit{kinēsis} is not merely a fading echo but a structured residue that retains the formal character of the original perception while being materially reconstituted in the organ's ongoing physiological activity. White's account of the twofold division of the soul's powers---\textit{kinēsis} on the one hand, \textit{aisthēsis} and \textit{gnōsis} on the other---provides a framework for understanding the resonant character of the persisting impression (White, \textit{The Meaning of Phantasia in Aristotle's De Anima III, 3-8}, p.~8). The residual \textit{kinēsis} is not a third power of the soul but a modality of the perceptual power itself; it is the perceptual capacity in its after-state, continuing to exercise a diminished form of the discriminative activity that characterizes it in full actuality. This is why Aristotle can say that imagination takes place in the absence of both the faculty and the activity of sense---``imagination takes place in the absence of both, as e.g. in dreams'' (Aristotle, \textit{On The Soul (De Anima)}, p.~42)---without contradicting his claim that \textit{phantasia} is causally derived from perception. The persisting impression is the bridge between the two: it is the residue of perception that enables imagination to occur in the absence of the original sensible object.

The signet-ring analogy, extended to three phases, illuminates what resonant \textit{kinēsis} adds to the hylomorphic account of §3. \textit{Phase one: potentiality.} The wax is pliable and smooth; the ring is held above. The sense-organ is ready and the sensible object nearby but unseen---$A_1$ is complete, $A_2$ not yet underway. \textit{Phase two: actualization.} The ring presses into the wax; the ring and wax are jointly actual in the impressing; ``the actual hearing and the actual sound come about at the same time,'' as Aristotle puts it at \textit{De Anima} III.2, 425b28--426a1. This is $A_2$ itself: the perceptual act as the shared actuality of sensible object and sense faculty. \textit{Phase three: residual trace.} The ring is withdrawn, the actualization proper ends---but the wax has undergone preservative \textit{alloiōsis}, and the impression persists. The \textit{aisthēmata} remain. The resonant \textit{kinēsis} is the organ's continuing exercise of the form received, sustained in the object's absence by the organ's own material substrate.

Burke's rhetorical analysis illuminates the resonant character of such persisting impressions in a different register. Discussing the strategies by which agents induce themselves and others to act, Burke notes that ``men induce themselves and others to act by devices that deduce `let us' from `we''' (Burke, \textit{A Grammar of Motives}, pp.~356--358). The rhetorical impression---the dispositional alteration produced by persuasive speech---similarly persists beyond the moment of its original delivery, shaping subsequent deliberation and action. The structural parallel with Aristotle's residual \textit{kinēsis} is perhaps instructive: in both cases, a temporally bounded act deposits a trace that continues to exert causal influence after the act itself has ceased. The resonance of the impression is not merely a metaphor but a description of the way in which the ensouled body---or, in Burke's terms, the motivated agent---retains and reactivates the formal structures deposited by past encounters. Von Uexküll's concept of \textit{Umwelt} provides yet another perspective on the persisting impression, one that emphasizes the organism's active constitution of its perceptual world. The \textit{Umwelt} is not a passive backdrop against which perception occurs but a meaningful environment that the organism constitutes through its perceptual and effectual organs (von Uexküll, \textit{A Foray Into the Worlds of Animals and Humans with A Theory of Meaning}, pp.~222--224). The resonant \textit{kinēsis} can thus be understood as a modification of the organism's \textit{Umwelt}: the persisting impression does not merely linger in the organ as a static trace but actively shapes the organism's subsequent encounters with its environment, determining what it notices, what it ignores, and how it responds.

The structural role of resonant \textit{kinēsis} in the chain follows from its character. An alteration that vanished with its cause could not ground further actualities; a completed actuality that terminated in itself could not serve as the unmoved originator of the next motion. But resonant \textit{kinēsis} persists, and in persisting, it becomes the material upon which the phantastic motion $M_2 \to A_3$ will operate. The \textit{phantasma} that will emerge at $A_3$ is not created \textit{ex nihilo} from the sensible object's form; it is generated from the resonant \textit{kinēsis} that $A_2$ has deposited in the organ. $A_2$'s residual trace is thus architecturally indispensable: without it, the chain could not continue past the perceptual event. The character of that residue, however---whether it is a single undifferentiated impression or a structured composite---determines what the subsequent stages of the chain can generate from it. The next subsection argues that the residue is not single but dual.

\subsection*{7. The Dual-Trace Thesis: Resonant Aisthēma and Resonant Orexis}

Here the section does its most consequential philosophical work. The persisting impression left by the completed perceptual act is not a single, undifferentiated residue but a dual trace comprising two distinct yet inseparable moments: the resonant \textit{aisthēma} (the persisting sensory form) and the resonant \textit{orexis} (the persisting appetitive valence). Every perception, in Aristotle's account, is simultaneously a discrimination of a sensible quality and an evaluation of that quality as pleasant or painful, beneficial or harmful; accordingly, every perception deposits in the ensouled body both a formal trace of what was perceived and an affective trace of how what was perceived was evaluated. This dual-trace thesis is perhaps the most consequential claim of the $A_1 \to A_2$ analysis, for it establishes the joint material and motivational ground from which \textit{phantasia}, memory, and action subsequently arise. The dual trace is numerically one event in the organ's substrate and two in being---the same formula Aristotle applies to motion at \textit{Physics} III.2, to perception at \textit{De Anima} III.2, and now, as this section argues, to the residual motion that perception leaves behind.

The argument proceeds along two converging lines. The first is analytic. Aristotle establishes at \textit{De Anima} II.3, 414b1--6 a chain of entailments: ``whatever has a sense has the capacity for pleasure and pain and therefore has pleasant and painful objects present to it, and wherever these are present, there is desire, for desire is appetition of what is pleasant.'' Every act of sensation, in this chain, analytically entails affective valence. To sense at all is already to be oriented toward the pleasant and away from the painful; to have pleasant and painful objects present is to have the appetite those objects engage. Sensation, on this account, is not a neutral registration of qualities that can then be evaluated by some further faculty; sensation is already evaluative in its very structure.

Aristotle develops the coordination of perception and appetite directly at \textit{De Anima} III.7, 431a8--14: ``To perceive then is like bare asserting or thinking; but when the object is pleasant or painful, the soul makes a sort of affirmation or negation, and pursues or avoids the object. To feel pleasure or pain is to act with the sensitive mean towards what is good or bad as such. Both avoidance and appetite when actual are identical with this: the faculty of appetite and avoidance are not different, either from one another or from the faculty of sense-perception; but their being is different.'' The passage is decisive. Perception and appetite are numerically one in substrate---one faculty, one event---but they differ in being. When the object is pleasant or painful, the perceptual act is simultaneously an affirmation or negation, a pursuit or avoidance. There is not first a perception and then, by a separate cognitive operation, an evaluation; the perception itself carries the evaluation as a formal aspect of its being.

The residual motion that persists after the perceptual act must inherit this dual structure. The sensation that deposited the trace was itself numerically one and two in being; its \textit{pathos} was formally both perceptual and orectic. Unless one can argue that the residue discards one aspect while retaining the other---which Aristotle's hylomorphism does not permit---the full \textit{pathos} persists, with both formal aspects intact. The enmattered-accounts argument at \textit{De Anima} I.1, 403a25--b19 supports this directly. The affections of soul are ``enmattered accounts'' (\textit{logoi enyloi}) whose definition must combine formal and material aspects; the physicist gives the material account (a boiling of the blood around the heart), the dialectician the formal account (a desire for retaliation), and the full definition requires both. If the residual motion is an enmattered \textit{pathos}, and the \textit{pathos} has two formally distinct aspects (perceptual and orectic), then the residue preserves both aspects as formal aspects of a single material substrate.

The second line of argument is a \textit{modus tollens}. If the residual motion were exhausted by its formal-epistemic aspect---if what persisted were only the resonant \textit{aisthēma}---then two consequences would follow that Aristotle's texts decisively reject. First, mere retention of sensible form would suffice to originate pursuit and avoidance. The soul's knowledge \textit{that the object is a wolf} would, by itself, move the animal to flee. But Aristotle insists that bare perception is ``like bare asserting or thinking''; it is only \textit{when the object is pleasant or painful} that the soul ``makes a sort of affirmation or negation, and pursues or avoids the object'' (\textit{De Anima} III.7, 431a8--14). The formal content alone is motivationally inert. A distinct operation---the hedonic-orectic valuation of the object as pleasant or painful---is required to initiate pursuit or avoidance. If the residue carried only the formal trace, it could not motivate; something more must persist.

Second, if only the formal trace persisted, the affective residue could have no independent operative force. But Aristotle observes at \textit{On Dreams} 2, 460b3--11 that ``when we are excited by emotions'' we are deceived by slight formal resemblances: ``the coward when excited by fear, the amorous person by the object of his desire; so that, with but little resemblance to go upon, the former thinks he sees his foes approaching, the latter, that he sees the object of his desire; and the more deeply one is under the influence of the emotion, the less similarity is required to give rise to these impressions.'' The passage establishes that affective residue primes formal recognition independently of the formal content. The coward and the courageous person, confronting the same minimal formal resemblance, respond differently---not because they have access to different formal traces but because their persisting affective states modulate the threshold at which formal recognition occurs. If the affective dimension were merely a shadow of the formal trace, no such differential priming would be possible; the coward and the courageous would recognize the same minimal resemblance or fail to recognize it identically. That they respond differently is evidence that the affective residue is doing its own operative work.

\textit{De Motu Animalium} 7, 702a17--19 provides the most architecturally explicit support: ``the organic parts are suitably prepared by the affections, these again by desire, and desire by imagination.'' Aristotle inserts the affections (\textit{pathē}) as a distinct causal step between desire and the organic parts. If the affective dimension were simply identical with the formal content of the \textit{phantasma}, its insertion as a distinct causal step would be otiose. Aristotle includes the \textit{pathē} because they do something the formal trace cannot do alone: they prepare the body for action by activating the hedonic-orectic axis. The \textit{pathē} at this point in the causal chain are the resonant \textit{orexis} at work---the affective residue that mediates between the phantasmatic representation and the bodily action. \textit{Rhetoric} I.11, 1370a27--b1 supplies complementary evidence in a different register: ``Pleasure is the consciousness through the senses of a certain kind of emotion; but imagination is a feeble sort of sensation, and there will always be in the mind of a man who remembers or expects something the imagination of what he remembers or expects. If this is so, it is clear that memory and expectation also, being accompanied by sensation, may be accompanied by pleasure.'' The pleasure that accompanies memory and expectation is not a fresh evaluation performed on a retrieved formal trace; it is a residual affection carried by \textit{phantasia} as part of the image's content. The image delivers the pleasure directly---because the pleasure was deposited alongside the formal trace in the original perceptual event, and because the dual-trace structure ensures that both persist together when the residue is activated.

Aristotle makes this duality explicit in his account of sensitive and deliberative imagination: ``Sensitive imagination, as we have said, is found in all animals, deliberative imagination only in those that are calculative: for whether this or that shall be enacted is already a task requiring calculation; and there must be a single standard to measure by, for that is pursued which is greater'' (Aristotle, \textit{On The Soul (De Anima)}, p.~52). The phrase ``that is pursued which is greater'' reveals the appetitive dimension inherent in even the most basic forms of \textit{phantasia}: the organism does not merely retain a sensory image but retains it under the aspect of its desirability, measuring it against alternatives. The Heideggerian observation that passions are ``both motion and a cause of motions''---a ``capacity for altering'' that ``alter[s] judgements,'' intruding upon \textit{logos} in a manner that reveals the inseparability of affect and cognition---offers a parallel formulation in a different idiom (Multiple Authors, \textit{Heidegger and Rhetoric}, pp.~113--115). The resonant \textit{orexis} is precisely this intrusion of affect upon form: it is the trace by which every perception carries not only information about the world but a dispositional orientation toward the world.

Burke's account of rhetorical motivation enriches this dual-trace analysis. In the \textit{Grammar of Motives}, Burke argues that the fiction of positive law serves to establish ``the values, traditions, and trends of business as the Constitution-behind-the-Constitution'' (Burke, \textit{A Grammar of Motives}, pp.~383--384), thereby revealing that every act of judgment presupposes a background of evaluative dispositions that are not themselves the objects of explicit deliberation. The resonant \textit{orexis} functions as a perceptual analogue of Burke's Constitution-behind-the-Constitution: it is the appetitive background against which every subsequent perceptual or cognitive act is evaluated, a background that is itself the product of prior perceptual encounters and their affective residues. Similarly, Burke's observation that ``antithesis is so strong a verbal instrument in both rhetoric and dialectic'' (Burke, \textit{A Rhetoric of Motives}, p.~203) illuminates the structural logic of the dual trace: the resonant \textit{aisthēma} and the resonant \textit{orexis} stand in a relation of complementarity rather than opposition, each requiring the other for the full significance of the perceptual event to be preserved. Frede's careful analysis of the cognitive role of \textit{phantasia} further supports the dual-trace thesis by noting that the ``uncontrolled status of \textit{phantasiai} as after-images'' does not do justice to their role ``in memory, thinking, or decision-making'' (Frede, \textit{The Cognitive Role of Phantasia in Aristotle}, pp.~4--5). The reason that \textit{phantasiai} can serve these multiple cognitive functions is precisely that they retain both formal and appetitive dimensions: a memory is not merely the recollection of a sensory image but the recollection of that image as pleasant or painful, significant or trivial; a deliberative \textit{phantasma} is not merely a picture of a possible future state but a picture charged with desire or aversion.

The convergence of the two arguments yields the dual-trace thesis in its final form. What persists after a perceptual event is not a single undifferentiated impression but a composite \textit{pathos} with two formally distinct aspects: the resonant \textit{aisthēma} (the ``what it is'' of the perceived object, the sensible form preserved) and the resonant \textit{orexis} (the ``how it matters for me'' of the object, the hedonic-orectic valence preserved). These two aspects are one in substrate---a single enmattered affection in the organ---but two in being, corresponding to the distinct ways of being of the perceptual and appetitive faculties that were already numerically one in the original act. Aristotle does not name the second trace with a technical term. He names the first as \textit{aisthēma} and treats affections as \textit{pathē} generally. The dual structure as developed here is a structural requirement of his own commitments---the commitment that sensation analytically entails affective valence (\textit{De Anima} II.3, 414b1--6), the commitment that perception and appetite are one in substrate and two in being (\textit{De Anima} III.7, 431a8--14), the commitment that the full \textit{pathos} of the sensory encounter persists after the object withdraws (\textit{On Dreams} 2, 459a24--28 and 460b2--3), and the commitment that the affective dimension has its own operative force within the causal chain of animal movement (\textit{De Motu Animalium} 7, 702a17--19). Take these commitments together, and the dual trace follows. The dual trace deposited by perception thus provides the material from which the entire subsequent cognitive life of the organism is constructed; it is the junction at which the physical world, as formally specified by the sensible object, is first translated into the motivational world of the ensouled body.

This account intersects with a productive tension in the corpus between \textit{phantasia} as causally derived from sensation and \textit{phantasia} as functionally autonomous in domains such as dreaming, rhetorical reception, memory retention, and involuntary movement. The dual-trace thesis offers a partial resolution of this tension: \textit{phantasia} is causally derived from perception insofar as the formal and appetitive traces that constitute it are deposited by the perceptual act, but it is functionally autonomous insofar as these traces, once deposited, can be reactivated, recombined, and transformed independently of the original sensible object. The resonant \textit{aisthēma} and the resonant \textit{orexis}, taken together, constitute the material basis for this functional autonomy; they are the seeds from which the virtual world of \textit{phantasia} grows, a world that is rooted in perception but not reducible to it.

\subsection*{8. Synthesis: $A_2$ as Ground of $A_3$}

The completed perceptual act, together with the dual residual trace it deposits, constitutes $A_2$---the second actuality of the perceptual chain and the ground upon which the next stage of actualization, $A_3$, will be built. $A_2$ is not merely an endpoint but a threshold: it names the state of the ensouled body after perception has occurred, a state in which the organism carries within itself the formal and appetitive determinations of its encounter with the sensible world. The section has traced $A_1 \to A_2$ through six moves: the preconditions at $A_1$ and the medium that links them; the hylomorphic analysis of form-reception; the unity of mover and moved as a single actuality under two \textit{logoi}; perception as preservative \textit{alloiōsis}; the persistence of residual motion as resonant \textit{kinēsis}; and the dual-trace thesis that specifies the internal structure of the residue. The transition from $A_2$ to $A_3$---from perception to \textit{phantasia}---will be the subject of the next section of this chapter; here it suffices to establish the precise sense in which $A_2$ serves as the ground of $A_3$.

$A_2$'s architectural role follows from the iterated chain structure established at $A_0$. Each completed actuality in the chain is simultaneously the terminus of its antecedent motion and the unmoved originator of its successor. $A_2$ is the terminus of the perceptual motion from $A_1$; it is also the unmoved originator of the phantastic motion $M_2 \to A_3$ through which the \textit{phantasma} proper will emerge. The resonant \textit{kinēsis} deposited at $A_2$ is the material upon which \textit{phantasia} will operate in the next stage. The \textit{phantasma} is not created \textit{ex nihilo} from the sensible object's form; it is generated from the dual trace that $A_2$ has left in the sensory apparatus. \textit{Phantasia}, as Aristotle defines it at \textit{De Anima} III.3, 428b10--17, is a \textit{kinēsis} that arises from actual sensation and is necessarily similar in character to the sensation from which it arose. The similarity is grounded in the resonant \textit{kinēsis}: the \textit{phantasma} inherits both the formal aspect (resonant \textit{aisthēma}, making the image recognizable as of that object) and the affective aspect (resonant \textit{orexis}, making the image matter to the perceiver) from the residual motion deposited at $A_2$.

First, then, $A_2$ supplies the formal content of \textit{phantasia}. Without the resonant \textit{aisthēma} deposited by perception, there would be no formal structure for \textit{phantasia} to preserve, recombine, or transform. Aristotle is explicit on this point: ``imagination takes place in the absence of both'' the faculty and the activity of sense (Aristotle, \textit{On The Soul (De Anima)}, p.~42), which means that imagination presupposes a prior perceptual act that has deposited a trace capable of being reactivated. Second, $A_2$ supplies the appetitive valence that charges \textit{phantasia} with motivational force. Without the resonant \textit{orexis}, the \textit{phantasma} would be a mere formal schema---a picture without significance---rather than the charged image that drives desire, aversion, and deliberation. Burke's analysis of how terminologies of motion and conditioning are dialectical enterprises designed to clarify the relationship between agent and scene applies here as well (Burke, \textit{A Grammar of Motives}, pp.~77--80); the dual trace deposited by perception ensures that both dimensions---formal and motivational---are preserved in the passage from $A_2$ to $A_3$.

The generalization established at $A_0$---each transition from one actuality to the next is itself a motion, a further actualization of a capacity that was latent in the preceding stage---applies here as it will at every subsequent stage. $A_2 \to A_3$ is the motion through which the residual trace is taken up by \textit{phantasia} and presented to the central organ as the \textit{phantasma} proper; $A_3 \to A_4$ is the motion through which the phantasmatic image, charged by the dual trace, mobilizes orexis toward action. Each transition is a motion in the \textit{energeia ateles} sense; each terminus is a completed actuality that becomes the unmoved originator of the next.

Furthermore, $A_2$ establishes the temporality of the actualization chain. Time, as the number of motion with respect to the before and after, provides the metric within which $A_0$, $A_1$, and $A_2$ are ordered; but $A_2$ introduces a new temporal dimension, namely the persistence of the trace beyond the moment of its deposition (Aristotle, \textit{Physics}, p.~60). The resonant \textit{kinēsis} is a temporally extended phenomenon: it is not merely a snapshot of the perceptual act but a continuing modification of the organism's psychophysical state that unfolds over time and is subject to decay, distortion, and reactivation. This temporal extension is what makes memory possible, and it is what distinguishes the Aristotelian account of perception from any model that treats the perceptual act as an instantaneous event without residual consequences. Hence $A_2$ is not merely the completion of one actualization but the inception of the next---the living hinge between perception and \textit{phantasia}, between the sensible world and the virtual world that the organism constructs from its perceptual residues. Gonzalez's emphasis on the reciprocal character of persuasive speech---the requirement that the orator's goodwill be reciprocated by the audience---provides a fitting analogy: just as the rhetorical act succeeds only when it leaves a persisting impression in the audience's soul, so too does the perceptual act succeed only when it deposits in the ensouled body a trace that grounds all subsequent cognitive and appetitive operations (Gonzalez, \textit{The Meaning and Function of Phantasia in Aristotle's `Rhetoric' III.1}, p.~29).

For now, the reader is deposited at the threshold of $A_3$, where the material deposited by perception becomes the object of the phantastic motion. The next section takes up that motion in detail, analysing how the resonant \textit{kinēsis} is activated by \textit{phantasia} and presented to the central organ, how the \textit{phantasma} that emerges is ontologically distinct from the trace on which it depends, and how the dual-aspect structure of the resonant \textit{kinēsis} is inherited and elaborated by the \textit{phantasma} that $A_3$ names.\footnote{The Heideggerian analysis of perception---particularly Heidegger's reading of \textit{aisthēsis} as a mode of Dasein's disclosive being-in-the-world and his treatment of the sensible object's presence (\textit{Anwesenheit}) as the ontological condition for perceptual encounter---illuminates what this section establishes. Heidegger's analysis of \textit{Befindlichkeit} (state-of-mind, attunement) as the existential structure through which Dasein is always already disposed toward its world offers an especially suggestive parallel to the dual-trace thesis developed in §7, insofar as \textit{Befindlichkeit} names the affective-orectic dimension of world-disclosure that cannot be reduced to the formal cognitive aspect. A full integration of the Aristotelian and Heideggerian accounts of perception lies beyond the scope of the present chapter and is developed in [Chapter X: The Motion-Inducing Soul---Rhetoric, Attunement, and Actualization (Heidegger, Rickert, Burke)].}

\subsection*{9. Development Note}

Three items flag future development. \textit{First}, the dual-trace thesis advanced in §7 treats the inheritance of affective valence as analytically entailed by sensation in all animals. This reading draws on \textit{De Anima} II.3, 414b1--6, which makes the entailment explicit for sensation as such. Whether every sensory affection in every animal inherits both traces, however, bears on the wider Aristotelian claim that appetite is the single source of all animal movement (\textit{De Anima} III.10, 433a27--31); and the question whether primitive organisms with only the sense of touch carry both traces or only the formal one merits engagement in a more extended treatment of animal locomotion. Nussbaum's commentary on the \textit{De Motu Animalium} and Caston's analysis of \textit{phantasia}'s role in animal cognition provide the principal points of reference for this debate.

\textit{Second}, the precise mechanism by which the medium's sequential affection translates into the unified perceptual act at $A_2$ is not fully articulated in \textit{De Anima} and requires supplementary analysis from \textit{Sense and Sensibilia} and the \textit{De Motu Animalium}. The sequential character of the medium's affection (\textit{Sense and Sensibilia} 6, 447a3--7) implies that the sensible form reaches the sense-organ over a temporal extension that varies by sensory modality, yet the perceptual act itself, at the moment of completion, presents the object as a unified formal whole rather than as a sequence of partial impressions. The mechanism through which the temporally-extended transmission issues in the temporally-unified perceptual experience is a genuine aporia that the Aristotelian corpus does not fully resolve; its treatment will draw on the \textit{koinē aisthēsis} analysis at \textit{De Anima} III.1--2 and on the further articulation of \textit{phantasia}'s synthesizing role in subsequent sections.

\textit{Third}, the dual-trace thesis developed here distinguishes the basic affective valence of resonant \textit{orexis}---the pursue-or-avoid, the pleasant-or-painful---from the determinate higher-order \textit{pathē} that Aristotle catalogues in the \textit{Rhetoric} (anger as desire accompanied by pain for conspicuous revenge; fear as disturbance at an anticipated destructive evil; pity as pain at an apparent evil befalling one who does not deserve it). The basic valence is co-constitutive with perception from $A_2$ onward and is inherited as a dispositional property by the \textit{phantasma} at $A_3$. The higher-order \textit{pathē}, by contrast, arise only downstream of \textit{doxa}'s propositional ratification at $M_2 \to M_3$---they require the rational cognitive apparatus and the committal dimension of opinion that this section's animals (which may include non-rational animals) do not necessarily possess. The distinction between \textit{pathos simpliciter} (basic valence, resonant \textit{orexis}) and \textit{pathē} in the higher-order sense (fear, anger, pity, shame) is structurally load-bearing for the downstream sections on cognitive actuality and emotion, where it will be developed in full.

\textit{Fourth}, the interpretive trajectory of the dissertation as a whole requires that each stage of the actualization chain be placed in dialogue with Heidegger's existential-phenomenological reappropriation of Aristotelian concepts. The Heideggerian register, sketched in preliminary form in the footnote to §8, will be developed at length in [Chapter X]. With respect to the opening subsection on the actualization of perception ($A_1 \to A_2$), Heidegger's treatment of \textit{aisthēsis} in his 1924 lecture on rhetoric recasts the Aristotelian account as a mode of Dasein's disclosive being-in-the-world, wherein perception is not a cognitive act performed by a subject upon an object but a fundamental manner in which being is disclosed (Multiple Authors, \textit{Heidegger and Rhetoric}, pp.~84--86). With respect to the preconditions ($A_1$ and the medium), Heidegger's concept of \textit{Befindlichkeit}---attunement or state-of-mind---re-reads Aristotle's structured readiness of the sense organ as an existential structure in which Dasein always already finds itself disposed toward a world; just as the Aristotelian organ must possess its first actuality as capacity, so too must Dasein be attuned before any particular being can be encountered. With respect to the unity of mover and moved, Heidegger's re-reading of Aristotle's \textit{Physics} III on \textit{kinēsis} re-interprets the shared actuality of perceiver and perceived as an ontological event in which beings are disclosed within the horizon of being; thus, the unity thesis is radicalized into the claim that being-in-the-world is always already a co-constitutive relation. With respect to perception as enmattered \textit{alloiōsis}, Heidegger observes that passions are ``both motion and a cause of motions''---a ``capacity for altering''---and that they ``alter judgements,'' thereby situating the Aristotelian doctrine of perceptual alteration within his own account of how \textit{Befindlichkeit} shapes \textit{logos} (Multiple Authors, \textit{Heidegger and Rhetoric}, pp.~113--115). With respect to resonant \textit{kinēsis}, Heidegger's analysis of the persisting character of mood---the way in which a past attunement continues to condition present disclosure---offers an existential analogue of Aristotle's residual movement. With respect to the dual-trace thesis, the Heideggerian parallel lies in the inseparability of understanding and state-of-mind: every disclosure of beings is simultaneously a cognitive determination and an affective disposition, a dual structure that mirrors the resonant \textit{aisthēma} and resonant \textit{orexis}. With respect to the synthesis of $A_2$ as ground of $A_3$, Heidegger's account of how past Dasein is retained in the present through historicality---the claim that historiological deficiency does not contradict but rather confirms Dasein's historicality---provides the existential framework for understanding how the perceptual trace grounds subsequent actualizations. Gibson's ecological approach to visual perception, which insists that ``the doctrine that vision is exteroceptive, that it obtains `external' information only, is simply false'' (Gibson, \textit{The Ecological Approach to Visual Perception}, pp.~209--210), will serve as a modern empirical counterpoint to both Aristotle's and Heidegger's accounts, bridging the ancient and phenomenological frameworks with a contemporary theory of perception as ecological engagement. Burke's rhetorical theory, particularly his account of identification, consubstantiality, and the scenic determinants of motivated action, will serve as the rhetorical counterpart that links perception, persuasion, and the formation of practical judgment across the entire actualization chain (Burke, \textit{A Rhetoric of Motives}, pp.~93--96).

\end{document}
